\textbf{Phần 1. Trắc nghiệm lựa chọn}
\setcounterpageref{ex}{0}

\begin{ex} %[3D2V71]
Tính giá trị của phép chia: $64\,824 : 4$.
\choice
{$16\,260$}
{$16\,204$}
{\True $16\,206$}
{$16\,106$}
\loigiai{
Chia lần lượt từ trái sang phải: $64\,824 : 4 = 16\,206$.
\textbf{Đáp án của câu hỏi là C}
}
\end{ex}

\begin{ex} %[3D2V71]
Một cuộn dây cáp dài $85\,450\text{ m}$ được chia đều thành 5 đoạn. Độ dài mỗi đoạn dây là
\choice
{\True $17\,090\text{ m}$}
{$17\,900\text{ m}$}
{$17\,080\text{ m}$}
{$16\,090\text{ m}$}
\loigiai{
$85\,450 : 5 = 17\,090\text{ (m)}$.
\textbf{Đáp án của câu hỏi là A}
}
\end{ex}

\begin{ex} %[3D2V71]
Tìm số tự nhiên $x$, biết: $x \times 6 = 78\,948$.
\choice
{$13\,148$}
{$13\,258$}
{$12\,158$}
{\True $13\,158$}
\loigiai{
$x = 78\,948 : 6 = 13\,158$.
\textbf{Đáp án của câu hỏi là D}
}
\end{ex}

\begin{ex} %[3D2V71]
Thực hiện phép chia có dư: $65\,487 : 5$. Thương và số dư lần lượt là
\choice
{$13\,097$ và dư $1$}
{\True $13\,097$ và dư $2$}
{$13\,970$ và dư $2$}
{$13\,096$ và dư $2$}
\loigiai{
$65\,487 = 5 \times 13\,097 + 2$ với số dư $2 < 5$.
\textbf{Đáp án của câu hỏi là B}
}
\end{ex}

\begin{ex} %[3D2V71]
Một kho hàng chia đều $48\,600\text{ kg}$ muối cho 6 cửa hàng. Mỗi cửa hàng đã bán được $\dfrac{1}{3}$ số muối của mình. Mỗi cửa hàng còn lại bao nhiêu ki-lô-gam muối?
\choice
{$2700\text{ kg}$}
{$8100\text{ kg}$}
{\True $5400\text{ kg}$}
{$6000\text{ kg}$}
\loigiai{
Mỗi cửa hàng nhận: $48\,600 : 6 = 8100\text{ kg}$. Đã bán: $8100 : 3 = 2700\text{ kg}$. Còn lại: $8100 - 2700 = 5400\text{ (kg)}$.
\textbf{Đáp án của câu hỏi là C}
}
\end{ex}

\begin{ex} %[3D2V71]
Trong một phép chia cho 7 có thương bằng $12\,450$ và số dư là số dư lớn nhất có thể có. Số bị chia là
\choice
{\True $87\,156$}
{$87\,150$}
{$87\,155$}
{$87\,157$}
\loigiai{
Số dư lớn nhất là 6. Số bị chia: $7 \times 12\,450 + 6 = 87\,150 + 6 = 87\,156$.
\textbf{Đáp án của câu hỏi là A}
}
\end{ex}

\begin{ex} %[3D2V71]
Một mảnh đất hình vuông có chu vi $65\,200\text{ m}$. Độ dài mỗi cạnh của mảnh đất đó là
\choice
{$16\,200\text{ m}$}
{$16\,400\text{ m}$}
{$15\,300\text{ m}$}
{\True $16\,300\text{ m}$}
\loigiai{
Cạnh hình vuông: $65\,200 : 4 = 16\,300\text{ (m)}$.
\textbf{Đáp án của câu hỏi là D}
}
\end{ex}

\begin{ex} %[3D2V71]
Tính giá trị của biểu thức: $A = 92\,480 : 8 - 4500$.
\choice
{$7160$}
{\True $7060$}
{$6960$}
{$7050$}
\loigiai{
$92\,480 : 8 = 11\,560$; $11\,560 - 4500 = 7060$.
\textbf{Đáp án của câu hỏi là B}
}
\end{ex}

\begin{ex} %[3D2V71]
Người ta chia đều $56\,700\ell$ xăng vào 7 bồn chứa. Mỗi bồn chứa bao nhiêu lít xăng?
\choice
{$8000\ell$}
{$8200\ell$}
{\True $8100\ell$}
{$7100\ell$}
\loigiai{
$56\,700 : 7 = 8100\text{ (\ell)}$.
\textbf{Đáp án của câu hỏi là C}
}
\end{ex}

\begin{ex} %[3D2V71]
Tìm một số biết rằng nếu lấy số đó chia cho 9 rồi bớt đi 1000 thì được 9000. Số đó là
\choice
{\True $90\,000$}
{$81\,000$}
{$80\,000$}
{$99\,000$}
\loigiai{
Số đó chia cho 9 bằng: $9000 + 1000 = 10\,000$. Số cần tìm: $10\,000 \times 9 = 90\,000$.
\textbf{Đáp án của câu hỏi là A}
}
\end{ex}

\begin{ex} %[3D2V71]
Một xí nghiệp có $72\,800$ chiếc áo cần đóng gói vào các thùng 8 chiếc. Số thùng đóng được là
\choice
{$9000$ thùng}
{$9200$ thùng}
{$9050$ thùng}
{\True $9100$ thùng}
\loigiai{
$72\,800 : 8 = 9100$ thùng.
\textbf{Đáp án của câu hỏi là D}
}
\end{ex}

\begin{ex} %[3D2V71]
Thực hiện phép chia: $100\,000 : 5$. Kết quả là
\choice
{$25\,000$}
{\True $20\,000$}
{$10\,000$}
{$50\,000$}
\loigiai{
$100\,000 : 5 = 20\,000$.
\textbf{Đáp án của câu hỏi là B}
}
\end{ex}

\textbf{Phần 2. Trắc nghiệm đúng sai}
\setcounterpageref{ex}{0}

\begin{ex} %[3D2V71]
Xét tính đúng sai của phép chia số có 5 chữ số cho số có 1 chữ số:
\choiceTF
{\True Phép chia thực hiện lần lượt từ trái sang phải.}
{\True $80\,040 : 8 = 10\,005$.}
{$80\,040 : 8 = 1005$.}
{Trong phép chia cho 6, số dư lớn nhất có thể là 6.}
\loigiai{
\begin{itemize}
    \item Ý a) Đúng: Quy tắc chuẩn.
    \item Ý b) Đúng: Chia có các chữ số 0 ở thương.
    \item Ý c) Sai: Sai vì thiếu chữ số 0 ở hàng chục nghìn/nghìn của thương.
    \item Ý d) Sai: Sai vì số dư lớn nhất là 5.
\end{itemize}
\begin{itemize}
    \item a) \textbf{Đúng}.
    \item b) \textbf{Đúng}.
    \item c) \textbf{Sai}.
    \item d) \textbf{Sai}.
\end{itemize}
}
\end{ex}

\begin{ex} %[3D2V71]
Một nông trường thu hoạch được $75\,000\text{ kg}$ cam chia đều cho 5 đội công nhân:
\choiceTF
{\True Mỗi đội nhận được $15\,000\text{ kg}$ cam.}
{\True Ba đội nhận được tất cả $45\,000\text{ kg}$ cam.}
{Hai đội nhận được nhiều hơn một nửa tổng số cam.}
{\True Tổng số cam chia hết, không dư ki-lô-gam nào.}
\loigiai{
\begin{itemize}
    \item Ý a) Đúng: $75\,000 : 5 = 15\,000\text{ kg}$.
    \item Ý b) Đúng: $15\,000 \times 3 = 45\,000\text{ kg}$.
    \item Ý c) Sai: Sai vì hai đội nhận 30 000 kg, ít hơn nửa (37 500 kg).
    \item Ý d) Đúng: 75 000 chia hết cho 5.
\end{itemize}
\begin{itemize}
    \item a) \textbf{Đúng}.
    \item b) \textbf{Đúng}.
    \item c) \textbf{Sai}.
    \item d) \textbf{Đúng}.
\end{itemize}
}
\end{ex}

\textbf{Phần 3. Trả lời ngắn}
\setcounterpageref{ex}{0}

\begin{ex} %[3D2V71]
Tính kết quả phép chia: $96\,848 : 8$.
\shortans{12106}
\loigiai{
\textbf{Đáp số:} $12106$\\
$96\,848 : 8 = 12\,106$.
}
\end{ex}

\begin{ex} %[3D2V71]
Tìm số tự nhiên $x$, biết: $x \times 7 = 84\,735$.
\shortans{12105}
\loigiai{
\textbf{Đáp số:} $12105$\\
$x = 84\,735 : 7 = 12\,105$.
}
\end{ex}

\begin{ex} %[3D2V71]
Một cuộn dây dài $54\,600\text{ m}$ chia đều thành 6 đoạn. Mỗi đoạn dài bao nhiêu mét?
\shortans{9100}
\loigiai{
\textbf{Đáp số:} $9100$\\
$54\,600 : 6 = 9100\text{ (m)}$.
}
\end{ex}

\begin{ex} %[3D2V71]
Tìm số dư của phép chia: $85\,693 : 9$.
\shortans{4}
\loigiai{
\textbf{Đáp số:} $4$\\
$85\,693 = 9 \times 9521 + 4$. Số dư là 4.
}
\end{ex}

\begin{ex} %[3D2V71]
Một hình vuông có chu vi $72\,400\text{ m}$. Tính độ dài cạnh hình vuông theo mét.
\shortans{18100}
\loigiai{
\textbf{Đáp số:} $18100$\\
$72\,400 : 4 = 18\,100\text{ (m)}$.
}
\end{ex}

\begin{ex} %[3D2V71]
Chia đều $90\,000\text{ kg}$ gạo cho 9 cửa hàng. Hỏi mỗi cửa hàng nhận được bao nhiêu ki-lô-gam gạo?
\shortans{10000}
\loigiai{
\textbf{Đáp số:} $10000$\\
$90\,000 : 9 = 10\,000\text{ (kg)}$.
}
\end{ex}
